\chapter{实验设计与结果分析}
\label{cha:experiment-analysis}

\section{实验目标与研究问题}
\label{sec:experiment-objectives}

实验部分的目标不是仅给出若干静态结果，而是验证第\ref{cha:system-implementation}章所述系统实现能否形成一条完整、可复现且具备分析价值的联邦入侵检测链路。结合 `docs/experimental\_design.md` 与现有实验产物，本文主要围绕以下问题展开验证：其一，轻量化 CNN-LSTM 模型是否能够在当前车联网演示数据上完成基本的攻击识别任务；其二，引入动态节点选择、Top-K 稀疏化与定点量化后，系统通信开销能否显著下降；其三，在通信压缩和节点筛选存在的情况下，检测性能是否保持在可接受范围内；其四，在原始数据不出本地的前提下，系统是否能够通过隐私代理指标体现出相对基线更高的综合收益。

根据实验设计文档，本文将上述目标进一步对应为四个研究问题。首先，研究问题 RQ1 关注动态节点选择是否有助于降低联邦轮次训练时延；其次，RQ2 关注参数压缩机制是否能够有效减少总通信量；再次，RQ3 关注改进方案对 Accuracy、Recall 和 F1 值等检测性能的影响是否可接受；最后，RQ4 关注在数据本地保留与上传更新压缩同时存在时，隐私代理分数是否能够得到提升。由此可见，实验章节强调的是“检测性能—通信收益—训练收益—隐私代理”四维协同分析，而不是单一精度指标的横向比较。

需要说明的是，当前仓库中能够直接核验的对照结果主要包括改进方案与标准 FedAvg 基线方案。实验设计文档中虽然列出了集中式训练和 FedProx 等候选对照组，但对应结果尚未在现有仓库产物中完整沉淀。因此，本章分析将以可直接核验的两组报告为主，并把其余对照方案保留为后续正式论文阶段的扩展实验方向。

\section{实验环境与参数设置}
\label{sec:experiment-setup}

当前实验使用的数据基础来自项目内置的 `data/demo\_vehicular.csv`，并以 `VeReMi \& Car-Hacking` 作为数据集名称写入配置和报告文件。需要指出的是，这里的数据集名称更多体现研究目标与后续适配方向，而当前仓库内直接提供的是演示性复现实验数据，用于验证工程链路和指标采集能力。预处理、训练和测试过程均通过 `main.py` 的命令行入口执行，推荐命令示例已整理于附录\ref{app:repro-commands}。

从运行环境看，项目 README 中建议使用 Python 3.10 及以上版本，并通过虚拟环境安装依赖。默认配置文件 `configs/default.toml` 指定设备为 CPU，数据路径为 `data/`，输出目录为 `outputs/`，同时给出了客户端数量、全局轮数、Top-K 比例、量化位宽和节点选择权重等默认参数。由于当前仓库中的实验结果主要用于工程验证，本章不把具体硬件平台作为核心比较变量，而是将平台信息交由 `run\_metadata.json` 自动记录，从而避免不同临时运行环境对结论表述造成混淆。

为了和已有报告保持一致，本文重点分析两组可核验方案：一组为改进方案，即动态节点选择 + Top-K 稀疏化 + 8 位量化；另一组为基线方案，即全客户端参与且不压缩上传更新的标准 FedAvg 形式。两组方案的关键参数设置如表\ref{tab:experiment-settings}所示。

\begin{table}[htbp]
    \centering
    \caption{实验主要参数设置}
    \label{tab:experiment-settings}
    \begin{tabular}{p{3.0cm}p{4.0cm}p{4.0cm}}
        \toprule
        参数项 & 改进方案 & 基线方案 \\
        \midrule
        数据集标识 & VeReMi \& Car-Hacking & VeReMi \& Car-Hacking \\
        客户端数量 & 4 & 4 \\
        联邦轮数 & 3 & 3 \\
        客户端参与率 & 0.75 & 1.0 \\
        Top-K 比例 & 0.2 & 1.0 \\
        量化位宽 & 8 bit & 32 bit \\
        节点选择权重 & 0.4 / 0.3 / 0.3 & 0.4 / 0.3 / 0.3 \\
        聚合方式 & FedAvg/FedProx 框架下的加权聚合 & 标准 FedAvg 对照 \\
        结果目录 & `outputs/` & `outputs/baseline\_fedavg/` \\
        \bottomrule
    \end{tabular}
\end{table}

需要特别说明的一点是，基线方案虽然仍保留了节点评分与选择证据结构，但由于 `client\_fraction=1.0`，所有可用客户端都会参与训练，因此评分结果不会影响是否被选中；相较之下，改进方案的 `client\_fraction=0.75` 使节点选择机制真正参与到了训练调度中。这一差异决定了后续关于时延和通信收益的比较具有明确的工程含义。

\section{评价指标与对照方案}
\label{sec:experiment-metrics-baselines}

为保证分析口径一致，本文沿用第\ref{sec:related-metrics}节中给出的评价指标定义，并结合当前工程原型的特点，将指标划分为检测性能、训练效率、通信效率与隐私代理四类。检测性能指标包括 Accuracy、Precision、Recall、F1 值和 False Positive Rate，用于衡量模型识别攻击样本的能力；训练效率指标包括总训练时间和平均轮次耗时，用于衡量联邦训练过程的整体时间开销；通信效率指标包括总原始通信量、总压缩后通信量和总体通信压缩率，用于评估参数压缩带来的直接收益；隐私代理指标则用于辅助说明“原始数据不出本地”与“上传更新量减少”对隐私暴露面的综合影响。

本章采用的对照方式遵循“单变量尽量少、工程意义尽量强”的原则。基线方案保留相同的模型结构和基本聚合逻辑，但让所有客户端都参与训练，并关闭压缩机制；改进方案在此基础上叠加节点选择、Top-K 压缩与 8 位量化，从而突出“调度 + 压缩”联合优化的综合收益。这样设计的优点在于：比较双方共享相同的模型骨架和训练框架，便于将观察到的差异主要归因于客户端参与策略和参数上传方式的变化。

除现有两组结果外，实验设计文档还提出了集中式训练、FedProx 和不同压缩比例、多种客户端参与率等扩展实验思路，这些内容有助于后续正式论文形成更完整的实验矩阵。但在当前阶段，本文仅将其作为后续工作保留，不将其写作“已完成实验”。这种处理方式既满足学术写作中事实边界清晰的要求，也有利于后续在真实 VeReMi 或 Car-Hacking 数据集上继续补充对照结果。

\section{实验结果与分析}
\label{sec:experiment-results}

根据 `outputs/reports/federated\_training\_report.md`、`outputs/baseline\_fedavg/reports/federated\_training\_report.md` 以及两组 `evaluation.json` 文件，可以得到改进方案与基线方案的主要对比结果，如表\ref{tab:experiment-main-results}所示。

\begin{table}[htbp]
    \centering
    \caption{改进方案与基线方案的关键结果对比}
    \label{tab:experiment-main-results}
    \begin{tabular}{p{4.2cm}p{2.8cm}p{2.8cm}p{2.0cm}}
        \toprule
        指标 & 改进方案 & 基线方案 & 对比结论 \\
        \midrule
        总训练时间（ms） & 933.24 & 1063.74 & 下降约 12.27\% \\
        平均轮次耗时（ms） & 304.51 & 332.83 & 下降约 8.51\% \\
        总通信量（bytes） & 79848 & 404928 & 下降约 80.28\% \\
        报告内通信压缩率 & 0.7371 & 0.0000 & 压缩效果明显 \\
        平均隐私代理分数 & 0.9211 & 0.7000 & 改进方案更高 \\
        Accuracy & 0.60 & 0.60 & 保持一致 \\
        Precision & 0.60 & 0.60 & 保持一致 \\
        Recall & 1.00 & 1.00 & 保持一致 \\
        F1-score & 0.75 & 0.75 & 保持一致 \\
        推理延迟（ms） & 3.45 & 3.14 & 差异较小 \\
        False Positive Rate & 1.00 & 1.00 & 均偏高 \\
        样本量 & 5 & 5 & 当前仅为演示验证 \\
        \bottomrule
    \end{tabular}
\end{table}

从通信侧结果看，改进方案的优势最为明显。由于系统在上传前仅保留 Top-K 关键增量，并进一步使用 8 位量化表达更新值，最终总通信量由基线方案的 404928 字节降至 79848 字节，下降幅度约为 80.28\%。报告中给出的总体通信压缩率为 0.7371，也表明在单轮原始上传量作为参考的情况下，压缩机制能够显著减少模型更新传输开销。对于资源受限且链路波动明显的车联网场景而言，这一收益具有直接工程价值。

从训练侧结果看，改进方案的总训练时间和平均轮次耗时也均优于基线方案。其原因主要体现在两个方面：一方面，改进方案把客户端参与率控制在 0.75，每轮仅选择部分节点参与训练，减少了需要串行完成的本地训练任务数；另一方面，参数压缩降低了每轮更新的处理与统计开销。虽然平均本地训练耗时在报告中略高于基线方案，但系统整体训练时间依然下降，说明“减少参与客户端数 + 降低上传量”的综合收益超过了局部训练时间波动带来的影响。

从检测侧结果看，改进方案与基线方案在 Accuracy、Precision、Recall 和 F1-score 上保持一致，分别为 0.60、0.60、1.00 和 0.75。这意味着在当前演示实验条件下，引入节点选择与压缩机制并未导致这些核心检测指标下降。与此同时，推理延迟分别为 3.45 ms 和 3.14 ms，差异较小，说明改进方案主要影响训练过程和上传机制，而不会明显增加本地推理负担。因此，可以初步认为当前系统在演示条件下实现了“通信开销显著下降而检测效果基本保持”的目标。

不过，结果中也暴露出一个必须正视的问题，即两组方案的 False Positive Rate 均为 1.00。这与 `evaluation.json` 中样本量仅为 5 的事实相互对应，并结合混淆矩阵可知，当前验证集中两个负类样本均被误判为正类。换言之，虽然召回率达到 1.00，但系统在当前极小样本演示条件下对正常样本的区分能力并不理想。因此，本文更适合将现阶段实验结果解释为“工程链路和指标采集机制可用”，而不能据此得出系统在真实车联网环境中已经具备稳定高精度检测性能的强结论。

从隐私代理维度看，改进方案的平均分数为 0.9211，高于基线方案的 0.7000。考虑到当前代理指标由本地性基础分与压缩收益共同构成，这一结果说明在保持原始数据留在本地的同时，进一步压缩上传更新量可以在工程意义上减少暴露面。但如前所述，该分数仅用于辅助比较，并不等价于差分隐私、同态加密或安全多方计算等严格隐私证明。

\section{局限性与有效性讨论}
\label{sec:experiment-limitations}

尽管当前实验已经能够支撑系统原型的功能验证和初步对照分析，但其局限性也非常明显。首先，当前实验使用的是仓库内置演示数据，最终评估样本量仅为 5，显然不足以代表真实车联网数据环境中的分布复杂性与攻击多样性。其次，现有结果主要集中在改进方案与标准 FedAvg 基线方案之间的单次对照，尚未对不同随机种子、不同参与率、不同 Top-K 比例和不同量化位宽开展系统性多轮重复实验，因此难以对结果稳定性给出统计意义上的判断。

其次，虽然实验设计文档提出了集中式训练、FedProx、多组压缩参数和多组节点选择权重等扩展方案，但这些内容在当前仓库中还未形成完整的结果证据链。因此，正式论文阶段仍需在真实 VeReMi 或 Car-Hacking 数据集上重新组织实验矩阵，至少对关键方案重复运行多次，并汇报均值与标准差，以增强结论的可信度。同时，还应补充更大规模验证集或独立测试集，避免当前“验证集即测试依据”的局限。

从有效性角度看，当前实验仍然具有一定价值。其一，它证明了系统已经具备从数据预处理到联邦训练、从模型评估到证据报告生成的完整闭环；其二，它给出了压缩机制和节点筛选对通信与训练效率影响的直接量化结果；其三，它揭示了在小样本场景下误报率偏高这一问题，为后续正式实验和模型改进指明了方向。因此，本文将现阶段结果定位为“工程验证性结果”是合适且必要的。

\section{本章小结}
\label{sec:experiment-summary}

本章结合实验设计文档、改进方案报告、基线方案报告与评估结果，对系统的实验目标、参数设置、对照方式和结果表现进行了分析。结果表明，在当前演示实验条件下，引入动态节点选择、Top-K 稀疏化与定点量化后，系统总通信量下降约 80.28\%，总训练时间下降约 12.27\%，平均轮次耗时下降约 8.51\%，而 Accuracy、Precision、Recall 与 F1-score 均未发生明显退化。同时，实验也暴露出样本量过小和误报率偏高等问题，说明现阶段结论更适合支持工程原型的有效性论证，而不宜直接外推到真实大规模车联网场景。下一章将在此基础上对全文工作进行总结，并给出后续扩展方向。
